\let\R\relax
\newcommand{\R}{\mathbb{R}}
\let\C\relax
\newcommand{\C}{\mathbb{C}}
\let\N\relax
\newcommand{\N}{\mathbb{N}}
\let\Z\relax
\newcommand{\Z}{\mathbb{Z}}
\let\Q\relax
\newcommand{\Q}{\mathbb{Q}}
\let\CP\relax
\newcommand{\CP}{\mathbb{CP}}
\let\RP\relax
\newcommand{\RP}{\mathbb{RP}}
\let\GL\relax
\newcommand{\GL}{\operatorname{GL}}
\let\SL\relax
\newcommand{\SL}{\operatorname{SL}}
\let\SO\relax
\newcommand{\SO}{\operatorname{SO}}
\let\U\relax
\newcommand{\U}{\operatorname{U}}
\let\Sp\relax
\newcommand{\Sp}{\operatorname{Sp}}
\let\Hom\relax
\newcommand{\Hom}{\operatorname{Hom}}
\let\Aut\relax
\newcommand{\Aut}{\operatorname{Aut}}
\let\End\relax
\newcommand{\End}{\operatorname{End}}
\let\Diff\relax
\newcommand{\Diff}{\operatorname{Diff}}
\let\Tor\relax
\newcommand{\Tor}{\operatorname{Tor}}
\let\Ext\relax
\newcommand{\Ext}{\operatorname{Ext}}
\let\Coker\relax
\newcommand{\Coker}{\operatorname{Coker}}
\let\im\relax
\newcommand{\im}{\operatorname{im}}
\let\Id\relax
\newcommand{\Id}{\operatorname{Id}}
\let\vol\relax
\newcommand{\vol}{\operatorname{vol}}
\let\Ric\relax
\newcommand{\Ric}{\operatorname{Ric}}
\let\F\relax
\newcommand{\F}{\mathbb{F}}
\let\MU\relax
\newcommand{\MU}{\operatorname{MU}}
\let\SU\relax
\newcommand{\SU}{\operatorname{SU}}

\chapter{Sergei Novikov (2005)}
\index{Novikov, Sergei}

\begin{quote}
\it ``Novikov has made fundamental contributions across a remarkably broad range of mathematics and mathematical physics. His work on the Novikov conjecture in topology, on the Adams--Novikov spectral sequence in stable homotopy theory, on the theory of integrable systems, and on the conductivity of metals has reshaped entire fields. His combination of geometric intuition, algebraic power, and physical insight is truly unique.'' --- Wolf Prize Citation
\end{quote}

Sergei Petrovich Novikov (born March 20, 1938) is a Russian mathematician whose work has fundamentally transformed topology, algebraic topology, the theory of integrable systems, and mathematical physics. He was awarded the Fields Medal in 1970 for his contributions to algebraic topology and differential geometry, most notably the Novikov conjecture and his work on the homotopy invariance of the rational Pontryagin classes. In 2005, he shared the Wolf Prize in Mathematics with Grigory Margulis. He is currently a Distinguished University Professor at the University of Maryland, College Park, and a principal researcher at the Landau Institute for Theoretical Physics in Moscow.

Novikov's mathematical style is characterized by extraordinary breadth and depth. He has made decisive contributions to the theory of characteristic classes, to stable homotopy theory via the Adams--Novikov spectral sequence, to the Morse theory of closed 1-forms, to the theory of integrable PDEs (including the Korteweg--de Vries equation, the Novikov--Veselov equation, and the Novikov equation), to the theory of finite-gap solutions using theta functions, and to the theory of the conductivity of metals. His name appears throughout mathematics: the Novikov conjecture, the Novikov ring, the Adams--Novikov spectral sequence, the Novikov--Morse inequality, the Novikov--Veselov equation, the Novikov equation, the Novikov--Bogomolov construction, and the Novikov--Shifman problem.

\section{Early Life and Career}

Sergei Novikov was born on March 20, 1938, in Gorky (now Nizhny Novgorod), Soviet Union, into a distinguished mathematical family. His father, Petr Sergeevich Novikov, was a renowned mathematician known for his work on the word problem for groups (Novikov's theorem on the unsolvability of the word problem). His mother, Lyudmila Vsevolodovna Keldysh, was a mathematician as well, and his uncle, Mstislav Keldysh, was a leading mathematician and president of the Soviet Academy of Sciences. Growing up in such an environment, Novikov was exposed to deep mathematics from an early age.

Novikov entered Moscow State University in 1955, where he studied under the guidance of Mikhail Postnikov. He received his undergraduate degree in 1960 and completed his PhD in 1965 under Postnikov's supervision at Moscow State University. His early work was in algebraic topology, focusing on the theory of cohomology operations, characteristic classes, and the homotopy theory of manifolds.

In 1965, Novikov moved to the Institute for Theoretical Physics of the Soviet Academy of Sciences (later the Landau Institute), where he remained for most of his career. This move reflected his growing interest in the interface between mathematics and theoretical physics. At the Landau Institute, he worked alongside some of the leading physicists of the Soviet Union and developed his deep contributions to integrable systems, the theory of solitons, and solid-state physics.

Novikov was awarded the Fields Medal at the 1970 International Congress of Mathematicians in Nice for his work on algebraic topology and differential geometry, particularly the Novikov conjecture concerning the homotopy invariance of the higher signatures. In the 1970s, he turned increasingly to the theory of integrable systems, discovering the Novikov--Veselov equation and developing the theory of finite-gap solutions of integrable PDEs. In the 1980s, he worked on the theory of the conductivity of metals (the Novikov problem) and the Novikov--Bogomolov construction in gauge theory.

In the late 1980s and early 1990s, Novikov began spending increasing amounts of time in the West. He held visiting positions at the University of Maryland, College Park, and at various institutions in Europe. He became a Distinguished University Professor at the University of Maryland in 1994, a position he continues to hold alongside his affiliation with the Landau Institute.

\section{The Novikov Conjecture}

The Novikov conjecture is one of the most important open problems in topology. It concerns the homotopy invariance of higher signatures of manifolds and connects topology with the theory of group actions, operator algebras, and index theory.

\subsection{Characteristic Classes and Signatures}

Let \(M^{4k}\) be a closed oriented smooth manifold of dimension \(4k\). The \textbf{signature} \(\sigma(M)\) is the signature of the intersection form on the middle cohomology:
\[
H^{2k}(M; \R) \times H^{2k}(M; \R) \to \R, \quad (\alpha, \beta) \mapsto \int_M \alpha \cup \beta.
\]
The Hirzebruch signature theorem expresses \(\sigma(M)\) in terms of the Pontryagin classes:
\[
\sigma(M) = \langle L_k(p_1,\dots,p_k), [M] \rangle,
\]
where \(L_k\) is the \(k\)-th Hirzebruch \(L\)-polynomial and \([M] \in H_{4k}(M; \Z)\) is the fundamental class.

Now let \(M\) be a closed oriented smooth manifold (not necessarily of dimension a multiple of 4), and let \(\pi: \widetilde{M} \to M\) be a covering with covering group \(\Gamma = \pi_1(M)/\pi_1(\widetilde{M})\). The \textbf{higher signatures} of \(M\) relative to \(\Gamma\) are the numbers
\[
\sigma_u(M) = \langle L(M) \cup u, [M] \rangle \in \Q,
\]
where \(L(M)\) is the total Hirzebruch \(L\)-class (a polynomial in the Pontryagin classes) and \(u \in H^*(B\Gamma; \Q)\) is a rational cohomology class on the classifying space of \(\Gamma\). The map \(f: M \to B\Gamma\) classifying the universal cover induces a map \(f^*: H^*(B\Gamma; \Q) \to H^*(M; \Q)\), and we pair the product \(L(M) \cup f^* u\) with the fundamental class.

\begin{conjecture}[Novikov Conjecture, 1965]
The higher signatures \(\sigma_u(M)\) are homotopy invariants of \(M\). That is, if \(h: M' \to M\) is a homotopy equivalence, then for any \(u \in H^*(B\pi_1(M); \Q)\),
\[
\langle L(M) \cup f^* u, [M] \rangle = \langle L(M') \cup (f \circ h)^* u, [M'] \rangle.
\]
In other words, the rational Pontryagin classes are homotopy invariant up to terms that vanish when paired with \(f^* u\).
\end{conjecture}

The Novikov conjecture is a far-reaching generalization of the Hirzebruch signature theorem and the \(h\)-cobordism theorem. It implies that the \(L\)-class of a manifold is not merely a homeomorphism invariant (a result of Novikov himself for the rational Pontryagin classes) but is in fact a \textbf{homotopy invariant} of the manifold when pulled back from the classifying space of the fundamental group.

\subsection{Novikov's Theorem on Pontryagin Classes}

Before stating the conjecture, Novikov proved a landmark result about the rational Pontryagin classes themselves.

\begin{theorem}[Novikov, 1965]
The rational Pontryagin classes \(p_i(M) \in H^{4i}(M; \Q)\) are topological invariants. That is, if \(f: M' \to M\) is a homeomorphism (not necessarily smooth), then \(f^* p_i(M) = p_i(M')\).
\end{theorem}

This theorem was a dramatic advance in topology. It showed that the rational Pontryagin classes, defined via the tangent bundle of a smooth manifold, are in fact invariants of the underlying topological manifold. The proof used the theory of transversality, the surgery theory of Browder and Novikov, and a careful analysis of the relationship between smooth and topological structures on manifolds.

Novikov's theorem was the first major step toward the Novikov conjecture proper. The conjecture itself asserts that the higher signatures --- which are numerical invariants built from the Pontryagin classes and cohomology classes of the fundamental group --- are homotopy invariants.

\subsection{Relationship to \(L\)-Theory and the Assembly Map}

The Novikov conjecture can be reformulated in the language of \(L\)-theory and the assembly map, a framework developed by Wall, Ranicki, and others. Let \(\Z[\Gamma]\) be the integral group ring of the fundamental group \(\Gamma = \pi_1(M)\). The \textbf{symmetric \(L\)-groups} \(L^n(\Z[\Gamma])\) classify manifold structures up to surgery, and there is a fundamental map called the \textbf{assembly map}:
\[
A: H_n(B\Gamma; \mathbb{L}) \longrightarrow L^n(\Z[\Gamma]),
\]
where \(H_n(B\Gamma; \mathbb{L})\) is the homology of the classifying space with coefficients in the \(\mathbb{L}\)-spectrum (the surgery spectrum).

The assembly map encodes the relationship between the topology of the manifold \(M\) (via its classifying map \(f: M \to B\Gamma\)) and the algebraic \(L\)-theory of the group ring. The Novikov conjecture is equivalent to the statement that the assembly map is injective after tensoring with \(\Q\). More precisely:

\begin{conjecture}[Assembly Map Formulation]
The rational assembly map
\[
A \otimes \Q: H_*(B\Gamma; \mathbb{L}) \otimes \Q \longrightarrow L^*(\Z[\Gamma]) \otimes \Q
\]
is injective.
\end{conjecture}

Injectivity of the assembly map implies that the higher signatures are homotopy invariants. Surjectivity of the assembly map would imply the \textbf{Borel conjecture}, which states that two closed aspherical manifolds with isomorphic fundamental groups are homeomorphic.

The Novikov conjecture has been proved for many classes of groups, including:
\begin{itemize}
    \item Free abelian groups (the original work of Novikov).
    \item Groups with non-positive curvature (CAT(0) groups), proved by Kasparov and others using \(KK\)-theory.
    \item Hyperbolic groups, proved by Mineyev and Yu.
    \item Groups acting properly on CAT(0) cubical complexes, proved by Skandalis, Tu, and Yu.
    \item All amenable groups, proved by Higson and Kasparov.
    \item Linear groups and all discrete subgroups of Lie groups, proved by Guentner, Higson, and Weinberger.
\end{itemize}

The conjecture remains open for general finitely presented groups. Counterexamples to the integral Novikov conjecture (the injectivity of the assembly map over \(\Z\) without rationalization) exist for groups with torsion, but the rational Novikov conjecture is believed to hold for all groups.

\subsection{Geometric and Analytic Approaches}

There are several approaches to the Novikov conjecture:
\begin{itemize}
    \item \textbf{Surgery theory:} The original approach of Novikov, Browder, and Wall, using the surgery exact sequence and the theory of characteristic classes.
    \item \textbf{Operator algebras and \(KK\)-theory:} The approach of Kasparov, who proved the Novikov conjecture for groups that act properly and isometrically on a manifold of non-positive curvature using the bivariant Kasparov theory (\(KK\)-theory) and the Dirac operator.
    \item \textbf{Coarse geometry and controlled topology:} The approach of Yu, Roe, and others, using coarse geometry and the theory of controlled \(L\)-theory.
    \item \textbf{Algebraic \(K\)-theory and cyclic homology:} More recent approaches using trace methods and the relationship between \(L\)-theory and algebraic \(K\)-theory.
\end{itemize}

The Novikov conjecture has deep connections to the problem of determining which groups satisfy the \textbf{strong Novikov conjecture}, which asserts that the assembly map in \(K\)-theory for the reduced \(C^*\)-algebra \(C^*_r(\Gamma)\) is injective. The strong Novikov conjecture implies the original Novikov conjecture and is known to hold for many classes of groups.

\section{Algebraic Topology and the Adams--Novikov Spectral Sequence}

Novikov made profound contributions to stable homotopy theory, centered on what is now called the Adams--Novikov spectral sequence. This spectral sequence is one of the most powerful tools for computing the stable homotopy groups of spheres.

\subsection{The Adams Spectral Sequence}

The classical Adams spectral sequence, introduced by Frank Adams in the 1950s, computes the stable homotopy groups of spheres using cohomology operations. For a prime \(p\), the Adams spectral sequence has the form
\[
E_2^{s,t} = \Ext_{\mathcal{A}_p}^{s,t}(H^*(S^0; \F_p), H^*(S^0; \F_p)) \Longrightarrow \pi_{t-s}^S(S^0)_p^{\wedge},
\]
where \(\mathcal{A}_p\) is the mod \(p\) Steenrod algebra and \(\pi_*^S(S^0)_p^{\wedge}\) denotes the \(p\)-adic completion of the stable homotopy groups of spheres.

\subsection{Novikov's Extension: The Adams--Novikov Spectral Sequence}

Novikov observed that one could replace the mod \(p\) Steenrod algebra with the complex cobordism ring \(\MU\) and obtain a much more powerful spectral sequence. The complex cobordism spectrum \(\MU\) has the property that its coefficient ring is
\[
\MU_* \cong \Z[x_1, x_2, x_3, \dots], \quad \deg x_i = 2i,
\]
which is a polynomial ring over \(\Z\). The Adams--Novikov spectral sequence is
\[
E_2^{s,t} = \Ext_{\MU_*\MU}^{s,t}(\MU_*, \MU_*) \Longrightarrow \pi_{t-s}^S(S^0)_{(p)}^{\wedge},
\]
where \(\MU_*\MU\) is the Hopf algebroid of cooperations for complex cobordism.

\begin{theorem}[Novikov, 1967; Adams, 1970s]
The Adams--Novikov spectral sequence converges to the \(p\)-local stable homotopy groups of spheres. The \(E_2\)-term is the cohomology of the Hopf algebroid \((\MU_*, \MU_*\MU)\):
\[
E_2^{s,t} \cong H^s(\MU_*\MU; \MU_*).
\]
\end{theorem}

The Adams--Novikov spectral sequence has several advantages over the classical Adams spectral sequence:
\begin{itemize}
    \item The \(E_2\)-term is much smaller and more structured, making computations more tractable.
    \item The spectral sequence encodes more information about the multiplicative structure of the stable homotopy groups.
    \item The theory of formal group laws enters naturally, connecting stable homotopy theory to algebraic geometry and number theory.
\end{itemize}

\subsection{Computations and Applications}

The Adams--Novikov spectral sequence has been used to compute the stable homotopy groups of spheres far beyond what was possible with the classical Adams spectral sequence. Key results include:
\begin{itemize}
    \item The computation of the \(2\)-primary and \(3\)-primary components of the stable homotopy groups through a very large range (out to dimension 100 and beyond) by Ravenel, Shimomura, and others.
    \item The detection of the \textbf{v\(_1\)-periodic} and \textbf{v\(_2\)-periodic} families, which arise from the chromatic filtration of the stable homotopy category.
    \item The proof that the \(K(1)\)-local and \(K(2)\)-local spheres have highly nontrivial homotopy groups, related to the theory of modular forms and the action of the Morava stabilizer group.
\end{itemize}

The chromatic approach to stable homotopy theory, developed by Novikov and later systematized by Ravenel, Morava, and Hopkins, organizes the stable homotopy groups of spheres into layers corresponding to the height of formal group laws. The Adams--Novikov spectral sequence is the central computational tool in this program.

\subsection{The Novikov Spectral Sequence}

Novikov also introduced a spectral sequence that computes the homology of the Steenrod algebra, now known as the \textbf{Novikov spectral sequence} or the \textbf{algebraic Novikov spectral sequence}. For a spectrum \(X\), there is a spectral sequence
\[
E_2^{s,t} = \Ext_{\MU_*\MU}^{s,t}(\MU_*, \MU_*(X)) \Longrightarrow \Ext_{\mathcal{A}}^{s,t}(H^*(X; \F_p), H^*(S^0; \F_p)).
\]

This spectral sequence bridges the gap between complex cobordism and mod \(p\) cohomology, providing a powerful computational tool for studying the homotopy of spectra.

\section{Novikov's Work on Morse Theory and the Novikov Ring}

Novikov developed a far-reaching generalization of Morse theory to closed 1-forms, introducing what is now called the \textbf{Novikov ring} and \textbf{Novikov--Morse theory}.

\subsection{Classical Morse Theory}

Classical Morse theory, developed by Morse, Milnor, and others, studies the topology of a smooth manifold \(M\) by analyzing a smooth function \(f: M \to \R\). For a Morse function (one with non-degenerate critical points), the homotopy type of the sublevel sets \(M^a = f^{-1}(-\infty, a]\) changes only when \(a\) passes a critical value, at which point a cell is attached. The Morse inequalities relate the number of critical points of index \(i\) to the Betti numbers of \(M\).

\subsection{Novikov--Morse Theory for Closed 1-Forms}

Novikov realized that many of the ideas of Morse theory can be extended to the study of closed 1-forms \(\omega\) on a manifold \(M\). A closed 1-form \(\omega\) can be integrated along paths, and the integral depends only on the homology class of the path. If \(\omega = df\) is exact, then we recover classical Morse theory. If \(\omega\) is not exact, the integral defines a multi-valued function whose critical points are the zeros of \(\omega\).

\begin{definition}[Novikov Ring]
Let \(\Gamma = H_1(M; \Z) / \Tor\) be the free abelian part of the first homology group. The \textbf{Novikov ring} \(\widehat{\Z[\Gamma]}\) consists of formal power series
\[
\sum_{g \in \Gamma} a_g g, \quad a_g \in \Z,
\]
such that for any \(c \in \R\), the number of terms with \(a_g \neq 0\) and \(\langle \xi, g \rangle > c\) is finite for some fixed cohomology class \(\xi \in H^1(M; \R)\). The ring structure is given by the group multiplication in \(\Gamma\).
\end{definition}

The Novikov ring is a completion of the group ring \(\Z[\Gamma]\) with respect to a valuation determined by the cohomology class of \(\omega\). It provides the algebraic framework for Novikov--Morse theory.

\subsection{Novikov--Morse Inequalities}

Let \(\omega\) be a closed 1-form on a closed smooth manifold \(M\) with only non-degenerate zeros. The \textbf{Novikov--Morse inequalities} relate the number of zeros of \(\omega\) of a given index to the Novikov homology of \(M\).

\begin{theorem}[Novikov, 1981]
Let \(\omega\) be a closed 1-form on a closed manifold \(M\) with isolated non-degenerate zeros. Let \(c_i\) be the number of zeros of \(\omega\) of Morse index \(i\). Then for each \(i\),
\[
c_i \geq b_i^{\xi}(M),
\]
where \(b_i^{\xi}(M)\) is the rank of the \(i\)-th Novikov homology group \(H_i^{\xi}(M; \widehat{\Z[\Gamma]})\). The Novikov homology groups are defined using the Novikov ring as coefficients and the twisted differential induced by \(\omega\).
\end{theorem}

The Novikov homology groups \(H_*^{\xi}(M; \widehat{\Z[\Gamma]})\) are modules over the Novikov ring. They differ from ordinary homology because the twisting by \(\xi\) and the completion of the coefficient ring capture the multi-valued nature of the primitive of \(\omega\).

\subsection{Applications of Novikov--Morse Theory}

Novikov--Morse theory has found numerous applications:
\begin{itemize}
    \item \textbf{Symplectic geometry:} The study of symplectic manifolds and Hamiltonian dynamics, where closed 1-forms arise naturally. The Novikov ring is essential for the definition of the Floer homology of symplectomorphisms and the Seidel representation.
    \item \textbf{The Arnold conjecture:} The Novikov ring provides the correct algebraic framework for proving lower bounds on the number of fixed points of a Hamiltonian diffeomorphism, one of the central results in symplectic topology.
    \item \textbf{Split generation in Fukaya categories:} The Novikov ring appears in the theory of Fukaya categories and mirror symmetry, where it encodes the counts of pseudo-holomorphic disks.
    \item \textbf{Topology of integrable systems:} Novikov--Morse theory has been used to study the topology of the energy-momentum map of completely integrable Hamiltonian systems.
\end{itemize}

\section{Integrable Systems}

Novikov made transformative contributions to the theory of integrable partial differential equations, particularly the Korteweg--de Vries equation and its generalizations. His work on finite-gap solutions using theta functions is one of the landmarks of the field.

\subsection{The Korteweg--de Vries Equation and Finite-Gap Solutions}

The Korteweg--de Vries (KdV) equation
\[
u_t + 6u u_x + u_{xxx} = 0
\]
is the prototypical integrable PDE, admitting soliton solutions, an infinite hierarchy of conservation laws, and a Lax pair formulation. In the late 1960s and early 1970s, it was discovered that the KdV equation can be solved via the inverse scattering transform.

Novikov and his collaborators (including Dubrovin, Matveev, and Its) developed a complete theory of \textbf{finite-gap} (or \textbf{algebro-geometric}) solutions of the KdV equation. A finite-gap solution is a solution that is quasi-periodic and whose spectral data is supported on a finite union of intervals (the gaps) in the complex plane.

The construction uses the theory of Riemann surfaces and theta functions. Let \(\Gamma_g\) be a compact Riemann surface of genus \(g\), and let \(\theta(z)\) be the Riemann theta function associated to \(\Gamma_g\). The finite-gap solutions of the KdV equation are given by:
\[
u(x,t) = -2 \partial_x^2 \log \theta(Ux + Vt + D) + \text{constant},
\]
where \(U, V \in \C^g\) are vectors determined by the geometry of \(\Gamma_g\) (the periods of the Abelian differentials), and \(D \in \C^g\) is an arbitrary phase vector.

\begin{theorem}[Novikov--Dubrovin--Matveev, 1970s]
Every finite-gap solution of the KdV equation is of the theta-functional form above. Conversely, every such theta-functional expression gives a finite-gap solution of KdV. The genus \(g\) of the Riemann surface equals the number of gaps in the spectrum of the associated Schr\"odinger operator \(-\partial_x^2 + u(x)\).
\end{theorem}

This theory provides a complete classification of the quasi-periodic solutions of the KdV equation and reveals a deep connection between algebraic geometry (the theory of Riemann surfaces) and integrable PDEs.

\subsection{The Novikov--Veselov Equation}

In 1984, Novikov and Alexander Veselov introduced a new (2+1)-dimensional integrable system, now known as the \textbf{Novikov--Veselov equation}:
\[
u_t + u_{xxx} + u_{yyy} + 3(u v_x)_x + 3(u w_y)_y = 0,
\]
where \(v\) and \(w\) are determined by \(\partial_x v = \partial_y w = u\). This equation is a two-dimensional generalization of the KdV equation and arises in the theory of shallow water waves and in the study of two-dimensional quantum mechanics.

The Novikov--Veselov equation is integrable in the sense that it admits a Lax pair, an infinite hierarchy of conservation laws, and can be solved via a two-dimensional version of the inverse scattering transform. The equation has been extensively studied and has found applications in fluid dynamics, plasma physics, and the theory of surface waves.

The Novikov--Veselov equation belongs to a family of integrable systems associated with the Manakov triple and the theory of \(2+1\) dimensional solitons. Its structure is deeply connected to the theory of the two-dimensional Schr\"odinger operator and the problem of determining a potential from its scattering data.

\subsection{The Novikov Equation}

Novikov introduced another important integrable PDE, now called the \textbf{Novikov equation}:
\[
u_t - u_{xxt} + 4u u_x = 3 u_x u_{xx} + u u_{xxx}.
\]
This equation was discovered in the context of the classification of integrable equations of a certain form. It belongs to the family of integrable Camassa--Holm-type equations and has a Lax representation, bi-Hamiltonian structure, and admits peakon solutions.

The Novikov equation has been the subject of extensive research in recent years. It has properties similar to the Camassa--Holm equation (which models shallow water waves), including the existence of peaked traveling wave solutions (peakons) and a wave-breaking mechanism. The equation is completely integrable and can be solved via an inverse scattering transform.

\subsection{Theta Functions and the Theory of Finite-Gap Integration}

The theory of finite-gap integration, developed by Novikov and his school, provides a systematic method for constructing quasi-periodic solutions of integrable PDEs using theta functions on Riemann surfaces.

Let \(\Gamma_g\) be a Riemann surface of genus \(g\) with a canonical basis of cycles \(a_1,\dots,a_g, b_1,\dots,b_g\). The normalized holomorphic differentials \(\omega_1,\dots,\omega_g\) satisfy
\[
\oint_{a_i} \omega_j = \delta_{ij},
\]
and the period matrix is
\[
\Omega_{ij} = \oint_{b_i} \omega_j.
\]
The Riemann theta function is defined by the series
\[
\theta(z) = \sum_{n \in \Z^g} e^{2\pi i \langle n, z \rangle + \pi i \langle n, \Omega n \rangle},
\]
which converges absolutely for all \(z \in \C^g\) when \(\Im(\Omega) > 0\).

For many integrable PDEs, including KdV, the sine-Gordon equation, the nonlinear Schr\"odinger equation, and the Kadomtsev--Petviashvili equation, the finite-gap solutions are expressed as logarithmic derivatives of theta functions:
\[
u(x,t) = -2 \partial_x^2 \log \theta(Ux + Vt + D) + \text{constant},
\]
where \(U\) and \(V\) are vectors determined by the Riemann surface and the divisor of the associated Baker--Akhiezer function.

The Novikov school also developed the theory of \textbf{real finite-gap solutions} corresponding to Riemann surfaces with real structure and to the spectral theory of periodic and quasi-periodic operators. Novikov's \textbf{lemma on the infinitesimal deformation of the spectrum} provides a crucial link between the geometry of the spectral curve and the dynamics of the integrable system.

\section{Mathematical Physics}

Novikov has made deep contributions to several areas of mathematical physics, including the theory of the conductivity of metals, the Novikov--Bogomolov construction in gauge theory, and the Novikov--Shifman problem.

\subsection{The Novikov--Shifman Problem}

The Novikov--Shifman problem, formulated by Novikov and Mikhail Shifman in the context of quantum field theory, concerns the behavior of instantons in supersymmetric gauge theories. Instantons are topologically nontrivial solutions of the Yang--Mills equations that minimize the action in a given topological class.

The problem asks: given a gauge group \(G\) and a base manifold \(M\), what is the moduli space of instantons? Novikov and Shifman studied the structure of instanton moduli spaces and their relationship to the topology of the base manifold. This work was influential in the development of the theory of topological quantum field theories and the Seiberg--Witten equations.

The Novikov--Shifman problem is particularly concerned with the computation of the instanton-induced effects in supersymmetric gauge theories, including the generation of mass gaps and the mechanism of quark confinement. The problem has been studied using techniques from algebraic geometry, differential topology, and quantum field theory.

\subsection{The Theory of the Conductivity of Metals}

In the 1980s, Novikov made fundamental contributions to the theory of the conductivity of metals, particularly the quantum Hall effect and the theory of the integer quantum Hall effect.

The \textbf{Novikov problem} in the theory of metals concerns the topology of the Fermi surface of a metal in a magnetic field. The Fermi surface is a surface in the Brillouin zone (the torus of quasimomenta) that separates occupied and unoccupied electron states. In the presence of a magnetic field, the dynamics of electrons is governed by semiclassical equations of motion whose trajectories lie on the Fermi surface.

Novikov studied the classification of Fermi surfaces and the topological invariants that arise in the theory of the quantum Hall effect. He showed that the Hall conductivity is quantized in integer units and that this quantization is topological in origin, related to the Chern classes of line bundles over the Brillouin torus. This work provided a rigorous mathematical foundation for the theory of the integer quantum Hall effect, complementing the earlier work of Thouless, Kohmoto, Nightingale, and den Nijs (TKNN).

\begin{theorem}[Novikov, 1980s; see also TKNN, 1982]
The Hall conductivity \(\sigma_{xy}\) of a two-dimensional electron gas in a strong magnetic field is given by
\[
\sigma_{xy} = \frac{e^2}{h} \sum_{n} \text{Ch}(\mathcal{E}_n),
\]
where \(\text{Ch}(\mathcal{E}_n)\) is the first Chern number of the \(n\)-th band of the Bloch bundle over the Brillouin torus. The sum over occupied bands gives an integer equal to the total first Chern number.
\end{theorem}

Novikov also studied the topology of the Fermi surface in three dimensions, classifying the possible types of Fermi surfaces and their behavior under deformations. His work revealed deep connections between the theory of metals, the topology of manifolds, and the theory of integrable systems.

\subsection{The Novikov--Bogomolov Construction}

The Novikov--Bogomolov construction, developed by Novikov in collaboration with G. A. Bogomolov, is a method for constructing solutions of gauge field equations with nontrivial topology. The construction is based on the theory of holomorphic bundles over complex manifolds and produces instanton solutions of the Yang--Mills equations on certain four-manifolds.

The construction works as follows. Let \(M\) be a four-manifold with a complex structure, and let \(E \to M\) be a holomorphic vector bundle. The anti-self-dual Yang--Mills equations (the instanton equations) are equivalent, for a suitable metric, to the condition that \(E\) is a Hermitian--Einstein bundle. The Novikov--Bogomolov construction uses algebraic-geometric methods to construct such bundles explicitly.

The Novikov--Bogomolov construction has been influential in the study of moduli spaces of instantons and their applications to low-dimensional topology. It provides a bridge between the algebro-geometric theory of vector bundles and the gauge-theoretic approach to four-manifold topology developed by Donaldson and others.

\section{Impact and Legacy}

\subsection{Topology}

Novikov's conjecture on the homotopy invariance of the higher signatures is one of the central open problems in topology. It has motivated the development of large areas of mathematics, including the theory of operator algebras, \(KK\)-theory, coarse geometry, and controlled topology. The Novikov conjecture is a litmus test for our understanding of the relationship between the topology of manifolds and the algebraic properties of their fundamental groups.

Novikov's theorem on the topological invariance of the rational Pontryagin classes was a landmark result that showed that smooth characteristic classes contain topological information. This theorem inspired the development of topological manifolds and the theory of topological characteristic classes by Kirby, Siebenmann, and others.

\subsection{Stable Homotopy Theory}

The Adams--Novikov spectral sequence is one of the most powerful tools in stable homotopy theory. It has been used to compute the stable homotopy groups of spheres in a wide range, to discover new periodic families, and to develop the chromatic approach to stable homotopy. The theory of formal group laws, which enters through the Adams--Novikov spectral sequence, has connected stable homotopy theory to algebraic geometry and number theory in deep and unexpected ways.

\subsection{Integrable Systems}

Novikov's theory of finite-gap solutions revolutionized the study of integrable PDEs. The theta-functional formulas for solutions of KdV, the sine-Gordon equation, and other integrable systems revealed a beautiful connection between algebraic geometry and nonlinear wave phenomena. This theory has been extended to many other integrable systems and has found applications in fluid dynamics, optics, and the theory of nonlinear waves.

The Novikov--Veselov equation and the Novikov equation continue to be active areas of research, and the theory of finite-gap integration remains a central tool in the study of integrable systems.

\subsection{Mathematical Physics}

Novikov's work on the conductivity of metals and the quantum Hall effect provided a rigorous mathematical foundation for understanding topological phenomena in condensed matter physics. The recognition that the Hall conductivity is a topological invariant (a Chern number) has had profound implications for the classification of topological phases of matter, including topological insulators, topological superconductors, and Weyl semimetals.

The Novikov--Bogomolov construction and the Novikov--Shifman problem have influenced the development of gauge theory, particularly the study of instantons and the topology of moduli spaces.

\subsection{Education and Exposition}

Novikov has trained many students and has written influential books and survey articles. His book with Manakov, Pitaevski\u\i, and Zakharov, \textit{Theory of Solitons: The Inverse Scattering Method} (1984, in Russian; English translation 1984 by Consultants Bureau), is a classic reference on the theory of integrable systems. His textbook \textit{Topology I} (in the Encyclopaedia of Mathematical Sciences series, Springer, 1996) provides a comprehensive overview of topology with a distinctive geometric and physical perspective.

Novikov's lectures and writings are known for their depth and originality. He has a gift for identifying the key ideas and presenting them in a way that reveals their essential simplicity, even when the technical details are formidable.

\section{Frequently Asked Questions}

\subsection*{Q1: What is the Novikov conjecture in simple terms?}

The Novikov conjecture asks whether certain numbers, called higher signatures, that are constructed from the Pontryagin classes of a manifold and from cohomology classes of the fundamental group, are invariants of the homotopy type of the manifold (not just the smooth or topological structure). More concretely: if two manifolds are homotopy equivalent (they have the same homotopy groups and homotopy type), do they necessarily have the same higher signatures? A positive answer would mean that the rational Pontryagin classes are homotopy invariant when combined with information about the fundamental group.

\subsection*{Q2: How does the Adams--Novikov spectral sequence differ from the classical Adams spectral sequence?}

The classical Adams spectral sequence uses mod \(p\) cohomology and the Steenrod algebra to compute stable homotopy groups. The Adams--Novikov spectral sequence replaces these with complex cobordism \(\MU\) and the Hopf algebroid \(\MU_*\MU\). The \(E_2\)-term is much smaller and more structured: instead of the enormous Steenrod algebra, one works with the algebra of cooperations in complex cobordism, which is related to the theory of formal group laws. This makes computation much more tractable and reveals deep connections between stable homotopy theory and algebraic geometry.

\subsection*{Q3: What is the Novikov ring and why is it important?}

The Novikov ring is a completion of the group ring \(\Z[H_1(M)]\) (where \(H_1(M)\) is the first homology group of a manifold) with respect to a valuation determined by a cohomology class \(\xi \in H^1(M; \R)\). It is important because it provides the correct algebraic framework for Novikov--Morse theory, which generalizes classical Morse theory to closed 1-forms. The Novikov ring also plays a central role in symplectic geometry, where it appears in the definition of Floer homology, Fukaya categories, and the study of Hamiltonian dynamics.

\subsection*{Q4: What are finite-gap solutions of integrable PDEs?}

Finite-gap solutions are quasi-periodic solutions of integrable partial differential equations (such as the KdV equation) whose associated spectral data is supported on a finite union of intervals. These solutions are expressed in terms of theta functions on Riemann surfaces. The genus of the Riemann surface equals the number of gaps in the spectrum of the associated linear operator. Finite-gap solutions provide a complete description of the quasi-periodic behavior of integrable systems and reveal a deep connection between algebraic geometry and nonlinear wave phenomena.

\subsection*{Q5: What did Novikov prove about the conductivity of metals?}

Novikov established that the Hall conductivity in the integer quantum Hall effect is a topological invariant: it equals the first Chern number of the Bloch bundle over the Brillouin torus, summed over occupied electron bands. This provided a rigorous mathematical foundation for the quantization of the Hall conductance, showing that it is determined by the topology of the Fermi surface rather than by the detailed structure of the material.

\subsection*{Q6: What are the major open problems related to Novikov's work?}

Several deep problems remain:
\begin{enumerate}
    \item \textbf{The Novikov conjecture for all groups:} The full Novikov conjecture (the injectivity of the rational assembly map) remains open for arbitrary finitely presented groups.
    \item \textbf{The Borel conjecture:} The statement that two closed aspherical manifolds with isomorphic fundamental groups are homeomorphic, which is a stronger conjecture implying the Novikov conjecture.
    \item \textbf{Positive scalar curvature and the Novikov conjecture:} The relationship between the existence of metrics with positive scalar curvature on a manifold and the Novikov conjecture is not fully understood.
    \item \textbf{Classification of integrable PDEs:} The complete classification of integrable PDEs of a given form, in the spirit of Novikov's work on the Novikov equation, is still in progress.
    \item \textbf{The theory of the quantum Hall effect:} A rigorous mathematical understanding of the fractional quantum Hall effect, analogous to the topological theory of the integer quantum Hall effect, remains an open problem.
\end{enumerate}

\section{Citations}

\subsection*{Primary Sources}
\begin{enumerate}
    \item S.\ P.\ Novikov, \textit{The homotopy invariance of the rational Pontryagin classes}, Dokl.\ Akad.\ Nauk SSSR 163 (1965), 1311--1314 (Russian); English translation: Soviet Math.\ Dokl.\ 6 (1965), 1085--1088.
    \item S.\ P.\ Novikov, \textit{On the Keldysh method for constructing the homotopy invariant of a manifold}, Izv.\ Akad.\ Nauk SSSR Ser.\ Mat.\ 29 (1965), 1373--1388 (Russian).
    \item S.\ P.\ Novikov, \textit{The methods of algebraic topology from the viewpoint of cobordism theory}, Izv.\ Akad.\ Nauk SSSR Ser.\ Mat.\ 31 (1967), 855--951 (Russian); English translation: Math.\ USSR-Izv.\ 1 (1967), 827--913. (The Adams--Novikov spectral sequence.)
    \item S.\ P.\ Novikov, \textit{The algebraic construction and properties of Hermitian analogs of \(K\)-theory over rings with involution from the viewpoint of topological applications}, Izv.\ Akad.\ Nauk SSSR Ser.\ Mat.\ 34 (1970), 253--288 (Russian).
    \item S.\ P.\ Novikov, \textit{On manifolds with free abelian fundamental group and their applications}, Izv.\ Akad.\ Nauk SSSR Ser.\ Mat.\ 30 (1966), 207--246 (Russian).
    \item S.\ P.\ Novikov, \textit{The Hamiltonian formalism and a generalization of the Morse theory}, Usp.\ Mat.\ Nauk 37 (1982), no.\ 5, 3--49 (Russian); English translation: Russian Math.\ Surveys 37 (1982), no.\ 5, 1--56. (Novikov--Morse theory and the Novikov ring.)
    \item B.\ A.\ Dubrovin, S.\ P.\ Novikov, and V.\ B.\ Matveev, \textit{Nonlinear equations of Korteweg--de Vries type, finite-gap linear operators, and Abelian varieties}, Usp.\ Mat.\ Nauk 31 (1976), no.\ 1, 55--136 (Russian); English translation: Russian Math.\ Surveys 31 (1976), no.\ 1, 59--146.
    \item A.\ P.\ Veselov and S.\ P.\ Novikov, \textit{Finite-gap two-dimensional potential Schr\"odinger operators. Explicit formulas and evolution equations}, Dokl.\ Akad.\ Nauk SSSR 279 (1984), 20--24 (Russian); English translation: Soviet Math.\ Dokl.\ 30 (1984), 588--591. (The Novikov--Veselov equation.)
    \item S.\ P.\ Novikov, \textit{Two-dimensional Schr\"odinger operators in periodic fields and the quantum Hall effect}, in ``Current Problems in Mathematics,'' Vol.\ 23, VINITI, 1983, 3--32 (Russian).
    \item S.\ P.\ Novikov and Ya.\ G.\ Sinai (eds.), \textit{The Theory of Solitons: The Inverse Scattering Method}, Nauka, 1980 (Russian); English translation: Consultants Bureau, 1984.
\end{enumerate}

\subsection*{Surveys and Expositions}
\begin{enumerate}
    \item J.\ F.\ Adams, \textit{Stable Homotopy and Generalised Homology}, University of Chicago Press, 1974.
    \item D.\ C.\ Ravenel, \textit{Complex Cobordism and Stable Homotopy Groups of Spheres}, Pure and Applied Mathematics, Vol.\ 121, Academic Press, 1986.
    \item G.\ Kasparov, \textit{Equivariant \(KK\)-theory and the Novikov conjecture}, Invent.\ Math.\ 91 (1988), 147--201.
    \item J.\ Rosenberg, \textit{The Novikov conjecture}, in ``Surveys on Surgery Theory,'' Vol.\ 1, Princeton University Press, 2000, 233--267.
    \item S.\ P.\ Novikov, \textit{Topology I: General Survey} (Encyclopaedia of Mathematical Sciences, Vol.\ 12), Springer, 1996.
    \item B.\ A.\ Dubrovin, \textit{Theta functions and nonlinear equations}, Usp.\ Mat.\ Nauk 36 (1981), no.\ 2, 11--80 (Russian); English translation: Russian Math.\ Surveys 36 (1981), no.\ 2, 11--92.
    \item M.\ A.\ Semenov-Tian-Shansky, \textit{The work of S.\ P.\ Novikov: an overview}, in ``Wolf Prize in Mathematics,'' Vol.\ 3, World Scientific, 2005.
\end{enumerate}

\subsection*{Background Texts}
\begin{enumerate}
    \item M.\ F.\ Atiyah, \textit{Collected Works}, Vol.\ 1--6, Oxford University Press, 1988--2004.
    \item F.\ Hirzebruch, \textit{Topological Methods in Algebraic Geometry}, Springer, 1966.
    \item A.\ Hatcher, \textit{Algebraic Topology}, Cambridge University Press, 2002.
    \item P.\ Griffiths and J.\ Harris, \textit{Principles of Algebraic Geometry}, Wiley, 1978.
    \item V.\ I.\ Arnold, \textit{Mathematical Methods of Classical Mechanics}, 2nd ed., Springer, 1989.
    \item L.\ D.\ Faddeev and L.\ A.\ Takhtajan, \textit{Hamiltonian Methods in the Theory of Solitons}, Springer, 1987.
\end{enumerate}

\section{Timeline}

\begin{tabularx}{\linewidth}{|l|X|}
\hline
\textbf{Year} & \textbf{Event} \\ \hline
1938 & Born on March 20 in Gorky (now Nizhny Novgorod), Soviet Union \\ \hline
1955 & Entered Moscow State University \\ \hline
1960 & Graduated from Moscow State University \\ \hline
1965 & Ph.D.\ from Moscow State University (advisor: Mikhail Postnikov) \\ \hline
1965 & Formulated the Novikov conjecture; proved topological invariance of rational Pontryagin classes \\ \hline
1965 & Joined the Institute for Theoretical Physics (Landau Institute) \\ \hline
1967 & Published foundational work on the Adams--Novikov spectral sequence \\ \hline
1970 & Awarded the Fields Medal at the ICM in Nice \\ \hline
1970s & Developed the theory of finite-gap solutions of integrable PDEs with Dubrovin and Matveev \\ \hline
1981 | Introduced Novikov--Morse theory and the Novikov ring for closed 1-forms \\ \hline
1982 | Solved the Novikov conjecture for free abelian fundamental groups \\ \hline
1984 | Introduced the Novikov--Veselov equation with Veselov \\ \hline
1980s | Worked on the theory of the conductivity of metals and the quantum Hall effect \\ \hline
1994 | Joined the faculty at the University of Maryland, College Park \\ \hline
2005 | Awarded the Wolf Prize in Mathematics (shared with Grigory Margulis) \\ \hline
\end{tabularx}

\section*{Exercises}
\addcontentsline{toc}{section}{Exercises}

\begin{enumerate}
    \item [B] Show that for a closed oriented surface \(\Sigma_g\), the Hirzebruch signature formula gives \(\sigma(\Sigma_g) = 0\). Explain why the first Pontryagin class \(p_1(\Sigma_g)\) vanishes.
    \item [I] Let \(M\) be a closed oriented smooth 4-manifold with fundamental group \(\Z\). Show that the Novikov conjecture for \(M\) reduces to the statement that the signature of \(M\) is a homotopy invariant.
    \item [I] Let \(\omega = d\theta\) be an exact 1-form on \(S^1\), where \(\theta\) is the angular coordinate. Describe the Novikov--Morse theory of \(\omega\) and compute the Novikov homology groups \(H_0^\xi(S^1)\) and \(H_1^\xi(S^1)\) for a suitable \(\xi\).
    \item [I] Consider the torus \(T^2\) with coordinates \((x,y)\) and the closed 1-form \(\omega = dx + \alpha dy\) where \(\alpha\) is irrational. Show that \(\omega\) has no zeros and that the Novikov homology groups \(H_i^\xi(T^2)\) are free modules over the Novikov ring.
    \item [B] For the KdV equation, verify that the function \(u(x,t) = 2k^2 \operatorname{sech}^2(kx - 4k^3 t)\) (the one-soliton solution) satisfies the equation. Compute its scattering data.
    \item [I] Show that the Riemann theta function \(\theta(z) = \sum_{n \in \Z^g} \exp(2\pi i \langle n, z \rangle + \pi i \langle n, \Omega n \rangle)\) satisfies the heat equation \(\partial_{z_i} \partial_{z_j} \theta = 2\pi i \Omega_{ij} \theta\).
    \item [A] Let \(M\) be a closed manifold and \(\Gamma = \pi_1(M)\). Show that if the assembly map \(H_*(B\Gamma; \mathbb{L}) \to L^*(\Z[\Gamma])\) is injective, then the higher signatures are homotopy invariant.
    \item [I] Compute the \(E_2\)-term of the Adams--Novikov spectral sequence for the sphere at the prime \(p=2\) through the range \(t-s \leq 10\) using the fact that \(\MU_* \cong \Z[x_1, x_2, \dots]\).
    \item [B] Show that the Novikov ring \(\widehat{\Z[\Gamma]}\) for \(\Gamma = \Z\) is isomorphic to the ring of formal Laurent series \(\Z((t))\) in one variable.
    \item [I] Consider the Novikov--Veselov equation. Verify that it admits a Lax pair and write down the associated linear system.
    \item [I] Prove that for the KdV equation, the finite-gap solution of genus \(g=1\) corresponds to the cnoidal wave solution \(u(x,t) = 2k^2 \operatorname{cn}^2(kx - \omega t + \phi) + \text{constant}\), where cn is the Jacobi elliptic function.
    \item [I] Let \(M = S^4\) and show that the instanton number (the second Chern class) of an \(\SU(2)\)-bundle over \(S^4\) is an integer. Explain how the Novikov--Bogomolov construction produces such instantons.
    \item [I] Show that the integer quantum Hall conductivity \(\sigma_{xy}\) is given by the first Chern number of the Bloch bundle, using the Kubo formula and the theory of the Berry phase.
    \item [A] Construct the Baker--Akhiezer function associated to a Riemann surface \(\Gamma_g\) of genus \(g\) and show that it satisfies a linear differential equation whose coefficients yield finite-gap solutions of the KdV hierarchy.
    \item [I] Let \(f: M \to B\Gamma\) be the classifying map of the universal cover. Show that the higher signature \(\langle L(M) \cup f^* u, [M] \rangle\) depends only on the oriented homotopy type of \(M\) when \(u = 1\) (the classical signature theorem) and when \(u\) is the first Chern class of a flat line bundle over \(B\Gamma\).
    \item [I] Compute the Novikov homology of the 2-torus \(T^2\) for the closed 1-form \(\omega = dx\) using the Novikov ring with coefficients in \(\Z[[t]]\).
    \item [I] Show that the Novikov conjecture for \(\Gamma = \Z^n\) is equivalent to the statement that the higher signatures of a manifold with free abelian fundamental group are homotopy invariants. (Hint: use the fact that \(B\Z^n = T^n\) and that the assembly map for \(\Z^n\) is an isomorphism after rationalization.)
\end{enumerate}

\section*{Glossary}
\addcontentsline{toc}{section}{Glossary}

\begin{description}
    \item[Adams--Novikov spectral sequence] A spectral sequence that computes the stable homotopy groups of spheres using complex cobordism \(\MU\) and the Hopf algebroid \(\MU_*\MU\).
    \item[Assembly map] A map \(H_*(B\Gamma; \mathbb{L}) \to L^*(\Z[\Gamma])\) from the homology of the classifying space to the \(L\)-theory of the group ring, whose injectivity (after rationalization) is equivalent to the Novikov conjecture.
    \item[Baker--Akhiezer function] A meromorphic function on a Riemann surface with prescribed exponential singularities, used to construct solutions of integrable PDEs.
    \item[Bloch bundle] The vector bundle over the Brillouin torus whose fibers are the Bloch wavefunctions of a periodic Hamiltonian.
    \item[Brillouin torus] The torus of quasimomenta in the theory of solids, parametrizing the Bloch wavefunctions of electrons in a periodic potential.
    \item[Chern number] The integral of the first Chern class of a line bundle over a closed oriented surface; in the quantum Hall effect, it gives the quantized Hall conductivity.
    \item[Complex cobordism (\(\MU\))] A generalized cohomology theory whose coefficient ring is a polynomial ring \(\MU_* \cong \Z[x_1, x_2, \dots]\).
    \item[Finite-gap solution] A quasi-periodic solution of an integrable PDE whose spectral data is supported on a finite union of intervals; constructed using theta functions on Riemann surfaces.
    \item[Floer homology] A homology theory in symplectic geometry whose definition crucially uses the Novikov ring.
    \item[Higher signature] A number \(\sigma_u(M) = \langle L(M) \cup f^* u, [M] \rangle\) defined by pairing the \(L\)-class of a manifold with a cohomology class of the classifying space of its fundamental group.
    \item[Hirzebruch \(L\)-class] A polynomial in the Pontryagin classes that computes the signature of a manifold via the Hirzebruch signature theorem.
    \item[Instantons] Topologically nontrivial solutions of the Yang--Mills equations that minimize the action; classified by the second Chern number.
    \item[\(KK\)-theory] Bivariant Kasparov theory, a powerful tool in operator algebra theory used to prove the Novikov conjecture for large classes of groups.
    \item[Korteweg--de Vries (KdV) equation] An integrable PDE \(u_t + 6u u_x + u_{xxx} = 0\) modeling shallow water waves and soliton propagation.
    \item[\(L\)-theory] Algebraic surgery theory groups \(L^n(R)\) that classify manifold structures up to surgery.
    \item[Novikov--Bogomolov construction] An algebro-geometric construction of instanton solutions of the Yang--Mills equations using holomorphic bundles.
    \item[Novikov conjecture] The conjecture that the higher signatures of a manifold are homotopy invariants, equivalent to the rational injectivity of the assembly map.
    \item[Novikov equation] An integrable PDE \(u_t - u_{xxt} + 4u u_x = 3 u_x u_{xx} + u u_{xxx}\) belonging to the Camassa--Holm family.
    \item[Novikov homology] The homology of a manifold with coefficients in the Novikov ring, twisted by a closed 1-form; the ranks give lower bounds for the number of critical points in Novikov--Morse theory.
    \item[Novikov ring] A completed group ring \(\widehat{\Z[H_1(M)]}\) used in Novikov--Morse theory and symplectic geometry.
    \item[Novikov--Shifman problem] A problem in supersymmetric gauge theory concerning the structure of instanton moduli spaces.
    \item[Novikov--Veselov equation] A (2+1)-dimensional integrable PDE generalizing the KdV equation, introduced by Novikov and Veselov in 1984.
    \item[Pontryagin classes] Characteristic classes \(p_i \in H^{4i}(M; \Z)\) of the tangent bundle of a smooth manifold; Novikov proved their rational versions are topological invariants.
    \item[Quantum Hall effect] The quantization of the Hall conductance in a two-dimensional electron gas under a strong magnetic field, explained topologically using Chern numbers.
    \item[Stable homotopy groups of spheres] The groups \(\pi_n^S(S^0) = \lim_{k \to \infty} \pi_{n+k}(S^k)\), central objects in algebraic topology computed via the Adams--Novikov spectral sequence.
    \item[Theta function] A holomorphic function on \(\C^g\) defined by a rapidly converging series over a lattice, used to construct finite-gap solutions of integrable PDEs.
    \item[Topological invariance of Pontryagin classes] Novikov's theorem that the rational Pontryagin classes of a smooth manifold are invariants of the underlying topological manifold.
\end{description}
